% IEEE conference paper — I2S TX AXI4-Lite IP on MicroBlaze V SoC
% !TeX spellcheck = en-US
% LTeX: language=en-US
% !TeX encoding = utf8
% !TeX program = lualatex
% !BIB program = bibtex

\documentclass[conference,a4paper,english]{IEEEtran}[2015/08/26]

\usepackage{balance}
\usepackage{iftex}
\usepackage{upquote}
\usepackage[main=english]{babel}
\usepackage[hyphens]{url}

\makeatletter
\g@addto@macro{\UrlBreaks}{\UrlOrds}
\makeatother

\ifluatex
  \usepackage[no-math]{fontspec}
  \usepackage{unicode-math}
  \usepackage{selnolig}
\else
  \usepackage[zerostyle=b,scaled=.75]{newtxtt}
  \usepackage[T1]{fontenc}
\fi

\usepackage[
  babel=true,
  expansion=alltext,
  protrusion=alltext-nott,
  final
]{microtype}
\DisableLigatures{encoding = T1, family = tt* }

\usepackage{graphicx}
\graphicspath{{../thesi2s/contents/figures/}}

\usepackage{amsmath}
\usepackage[dvipsnames, table]{xcolor}
\usepackage{listings}

\definecolor{eclipseKeywords}{RGB}{127,0,85}
\definecolor{eclipseStrings}{RGB}{42,0,255}

\lstset{
  basicstyle=\footnotesize\ttfamily,
  breaklines=true,
  breakatwhitespace=true,
  columns=flexible,
  tabsize=2,
  numbers=left,
  numberstyle=\tiny,
  xleftmargin=.5cm,
  captionpos=b,
  keywordstyle=\color{eclipseKeywords},
  stringstyle=\color{eclipseStrings},
  commentstyle=\itshape
}

\lstloadlanguages{[LaTeX]TeX, C, Verilog}

\usepackage{array}
\usepackage{booktabs}
\usepackage{paralist}
\usepackage[
  square,
  comma,
  numbers,
  sort&compress
]{natbib}
\renewcommand{\bibfont}{\normalfont\footnotesize}

\usepackage{etoolbox}
\makeatletter
\patchcmd{\NAT@test}{\else \NAT@nm}{\else \NAT@hyper@{\NAT@nm}}{}{}
\makeatother

\usepackage{multirow}
\usepackage{stfloats}
\fnbelowfloat

\usepackage[group-minimum-digits=4,per-mode=fraction]{siunitx}
\usepackage{hyperref}
\hypersetup{
  hidelinks,
  colorlinks=true,
  raiselinks=true,
  allcolors=black,
  pdfstartview=Fit,
  breaklinks=true,
  hypertexnames=false,
}
\usepackage[all]{hypcap}
\usepackage[caption=false,font=footnotesize]{subfig}
\usepackage[capitalise,nameinlink,noabbrev]{cleveref}

\crefname{listing}{Listing}{Listings}
\Crefname{listing}{Listing}{Listings}
\crefname{lstlisting}{Listing}{Listings}
\Crefname{lstlisting}{Listing}{Listings}

\usepackage{xspace}
\newcommand{\eg}{e.g.,\ }
\newcommand{\ie}{i.e.,\ }

\input{ushyphex}
\hyphenation{Micro-Blaze Vivado}

\input{commands}

\ifpdftex
  \input{glyphtounicode}
  \pdfgentounicode=1
\fi

\begin{document}

\title{Design of an AXI4-Lite I2S Transmitter IP with\\DDS-Based Arbitrary Sample Rate for MicroBlaze V SoC}

\author{%
  \IEEEauthorblockN{Muhammad Muqtada Alhaddad}
  \IEEEauthorblockA{Department of Electrical Engineering and Information Technology\\
    Universitas Gadjah Mada\\
    Yogyakarta, Indonesia\\
    \texttt{22/500341/TK/54841}}
}

\maketitle

\begin{abstract}
This paper presents the design and implementation of a custom Inter-IC Sound (I2S) transmitter peripheral in Verilog HDL, integrated as an AXI4-Lite slave into a System-on-Chip (SoC) based on the MicroBlaze~V soft processor on a Xilinx Artix-7 FPGA.
The peripheral is controlled through memory-mapped registers from bare-metal software in Xilinx Vitis.
Audio clocks (BCLK, WS, and MCLK) are generated by a direct digital synthesis (DDS) phase accumulator running at 100\,MHz, enabling programmable arbitrary sample rates without reconfiguring external PLL blocks.
Shadow registers and a STATUS-bit handshake synchronize firmware sample delivery with hardware frame boundaries.
The system was validated on a Basys3 board with a PCM5102A DAC through oscilloscope measurements, synthesis reports, and listening tests.
Results show correct Philips I2S timing, sub-0.001\% sample-rate accuracy, and stable stereo audio output.
\end{abstract}

\begin{IEEEkeywords}
MicroBlaze V, I2S, Verilog, AXI4-Lite, FPGA, SoC, DDS, digital audio
\end{IEEEkeywords}

\IEEEpeerreviewmaketitle

\section{Introduction}
\label{sec:introduction}

Modern embedded systems increasingly require digital audio interfaces that are software-configurable and easy to integrate into system-on-chip (SoC) architectures.
The Inter-IC Sound (I\textsuperscript{2}S) protocol~\cite{philips_i2s_1986} is widely used to stream pulse-code modulation (PCM) samples from a processor to external DACs and codecs using three synchronous signals: bit clock (BCLK), word select (WS/LRCLK), and serial data (SD).
Field-programmable gate arrays (FPGAs)~\cite{trimberger_history_fpga_2018} enable custom real-time serializers and register interfaces in Verilog~\cite{palnitkar_verilog_2003}, while soft processors such as MicroBlaze~V~\cite{xilinx_microblaze_ref_2021,waterman_riscv_2016} provide a familiar software layer for configuration and test data generation through memory-mapped I/O (MMIO) over AXI4-Lite~\cite{arm_amba_axi_2010}.

Commercial I2S LogiCORE IP~\cite{amd_pg308_i2s_logicore} targets AXI4-Stream and multi-channel streaming use cases.
This work instead delivers a compact, TX-only peripheral with an AXI4-Lite register file, transparent RTL, and programmable sample rate via direct digital synthesis (DDS)~\cite{vankka_dds_2001}.
The design is vendor-agnostic at the audio-clock level because BCLK, WS, and MCLK are synthesized digitally from a single 100\,MHz system clock rather than from dedicated PLL audio outputs.

The contributions of this paper are:
\begin{compactitem}
  \item a custom I2S transmitter IP with AXI4-Lite control, shadow registers, and STATUS-bit hardware--software handshake;
  \item DDS-based arbitrary sample-rate generation up to 195.312\,kHz on a 100\,MHz clock domain;
  \item end-to-end integration on a Basys3 Artix-7 board~\cite{xilinx_basys3_ref} with MicroBlaze~V, bare-metal Vitis firmware~\cite{xilinx_vitis_embedded_2021}, and PCM5102A DAC validation~\cite{pcm5102a_datasheet}; and
  \item documented firmware defects and their fixes using closed-loop STATUS polling instead of \texttt{usleep()}-based timing.
\end{compactitem}

The remainder of the paper is organized as follows.
\Cref{sec:related} reviews related work.
\Cref{sec:architecture} describes the proposed architecture.
\Cref{sec:implementation} summarizes RTL and software implementation.
\Cref{sec:results} presents measurement results.
\Cref{sec:conclusion} concludes the paper.

\section{Related Work}
\label{sec:related}

Prior FPGA I2S implementations demonstrate low-latency audio transport at rates up to 96\,kHz~\cite{sang_i2s_codec_2020} and instrumented interfaces on low-cost boards~\cite{kaviani_i2s_fpga_2017}.
Most published designs focus on correct waveform generation in isolated RTL blocks.
Fewer works document a complete path from custom peripheral through AXI SoC integration to interactive bare-metal software validation~\cite{lafosse_embedded_2018,sklyarov_axi_peripheral_2014,crockett_zynq_book_2014}.

Hierarchical AXI-based SoCs on FPGA show that custom slaves can be attached modularly through interconnect fabrics~\cite{sklyarov_axi_peripheral_2014}.
Clock-domain crossing (CDC) is critical in multi-clock audio systems~\cite{ginosar_metastability_2011}.
The final design presented here operates entirely in the 100\,MHz \texttt{S\_AXI\_ACLK} domain, eliminating CDC on the PCM data path while using a STATUS toggle for firmware synchronization.

Compared with vendor I2S cores that rely on FIFOs and AXI-Stream~\cite{amd_pg308_i2s_logicore}, this work emphasizes register-mapped programmed I/O, explicit shadow latching, and software-visible DDS rate control suitable for academic replication on a Basys3 platform.

\section{Proposed Architecture}
\label{sec:architecture}

\subsection{System Overview}

The target system couples a MicroBlaze~V processor, AXI interconnect, block RAM, UART, and the custom I2S transmitter IP on a Basys3 FPGA.
Audio samples are written by firmware to \texttt{DATA\_LEFT} and \texttt{DATA\_RIGHT}; the IP serializes them in Philips I\textsuperscript{2}S format to a PCM5102A module over PMOD wiring using only three signals: Bit Clock (BCLK), Word Select (WS/LRCK), and Data In (DIN).

\begin{figure}[!t]
  \centering
  \IfFileExists{processor-uart-timer-axi-i2s-architecture-whole-diagram.png}{%
    \includegraphics[width=\linewidth]{processor-uart-timer-axi-i2s-architecture-whole-diagram.png}%
  }{}
  \caption{SoC architecture with custom I2S TX IP and external PCM5102A DAC (using BCLK, WS, and DIN).}
  \label{fig:arch-soc}
\end{figure}

\subsection{Register Map and Control Fields}

Four 32-bit AXI4-Lite registers expose the peripheral to software.
\Cref{tab:regmap} lists offsets and functions; \texttt{CONTROL} packs \texttt{ENABLE}, \texttt{MUTE}, \texttt{SAMPLE\_WIDTH} (1--32 bits, 0 means 32), and \texttt{SAMPLE\_RATE\_HZ} (0--1{,}048{,}575\,Hz, safely clamped to 195.312\,kHz).

\begin{figure}[!t]
  \caption{AXI4-Lite register map of the I2S transmitter IP.}
  \label{tab:regmap}
  \centering
  \begin{tabular}{@{}llll@{}}
    \toprule
    \textbf{Offset} & \textbf{Name} & \textbf{Acc.} & \textbf{Function} \\
    \midrule
    \texttt{0x00} & \texttt{DATA\_LEFT}  & W & Left PCM sample \\
    \texttt{0x04} & \texttt{DATA\_RIGHT} & W & Right PCM sample \\
    \texttt{0x08} & \texttt{CONTROL}     & R/W & Enable, mute, width, rate \\
    \texttt{0x0C} & \texttt{STATUS}      & R & Frame-latch toggle (bit~0) \\
    \bottomrule
  \end{tabular}
\end{figure}

\subsection{DDS Clock Generation}

BCLK is derived from a phase accumulator running at $f_{\mathrm{ref}} = 100$\,MHz.
The toggle rate is $128 \times F_s$ because each BCLK half-period is generated by the accumulator wrap logic.
With a 32-bit phase path and modulus tied to 100{,}000{,}000, the average output frequency error stays below 0.001\% for standard and non-standard rates.
Although MCLK is also synthesized at $256 \times F_s$ for general compatibility, it was not required for the PCM5102A validation as the DAC was configured to generate its internal master clock from BCLK.

For stereo 32-bit slots (64 BCLK periods per frame), Philips timing requires WS transitions one BCLK before the next channel MSB and a mandatory one-bit delay after each WS edge~\cite{philips_i2s_1986}.

\subsection{Data Integrity and Handshake}

Samples written by the CPU are copied into internal shadow registers only at the start of a stereo frame (\texttt{bit\_count == 0} on a rising BCLK edge).
This prevents mid-frame corruption if software updates a data register while serialization is in progress.
\texttt{STATUS[0]} toggles whenever a new left/right pair is latched, enabling firmware to poll for the next transmit slot instead of relying on wall-clock delays or the RISC-V \texttt{mcycle} counter~\cite{riscv_isa_vol1_2019}.

\begin{figure}[!t]
  \centering
  \IfFileExists{i2s-ip-block.png}{%
    \includegraphics[width=\linewidth]{i2s-ip-block.png}%
  }{}
  \caption{Internal structure of the DDS-based I2S transmitter IP.}
  \label{fig:ip-block}
\end{figure}

\section{Implementation}
\label{sec:implementation}

\subsection{RTL Design}

The IP was implemented in Verilog and packaged for Vivado IP Integrator~\cite{xilinx_vivado_design_suite_2021}.
The AXI4-Lite slave decodes single 32-bit read/write transactions to the four registers.
DDS logic in \texttt{i2s\_dds\_slave\_lite\_v1\_0\_S00\_AXI.v} computes \texttt{bclk\_toggle\_rate\_hz = sample\_rate\_hz << 7} (i.e., $F_s \times 128$) and toggles BCLK when the phase sum reaches 100{,}000{,}000.

\begin{lstlisting}[language=Verilog, caption={DDS phase accumulator for BCLK generation.}, label={lst:dds}]
wire [31:0] bclk_toggle_rate_hz = {sample_rate_sync, 7'b0};
wire bclk_phase_wrap = (bclk_phase_sum >= AXI_CLK_HZ);

always @(posedge S_AXI_ACLK) begin
    if (bclk_phase_wrap) begin
        bclk_phase_acc_q <= bclk_phase_sum - AXI_CLK_HZ;
        bclk_q <= !bclk_q;
    end else begin
        bclk_phase_acc_q <= bclk_phase_sum;
    end
end
\end{lstlisting}

At each stereo frame boundary the serializer copies \texttt{slv\_reg0}/\texttt{slv\_reg1} into shadow registers and flips the latch-status bit exported through \texttt{STATUS}.

\subsection{SoC Integration and Constraints}

The block design connects MicroBlaze~V, AXI SmartConnect, local BRAM, Clocking Wizard (100\,MHz system clock), AXI UARTLite, and the custom I2S IP.
Pin constraints map \texttt{i2s\_mclk}, \texttt{i2s\_bclk}, \texttt{i2s\_ws}, and \texttt{i2s\_data} to PMOD pins for the PCM5102A module.
The hardware platform is exported to Vitis as an \texttt{.xsa} file for bare-metal application development.

\subsection{Firmware}

The test application (\texttt{helloworld.c}) configures \texttt{CONTROL}, generates sine or square test tones with a software DDS phase accumulator, and streams samples using STATUS polling:

\begin{lstlisting}[language=C, caption={STATUS-polling sample loop.}, label={lst:firmware}]
while (1) {
    while (frame_status_bit() == g_last_frame_bit) {
        if (!XUartLite_IsReceiveEmpty(STDIN_BASEADDRESS))
            handle_key(XUartLite_RecvByte(STDIN_BASEADDRESS));
    }
    g_last_frame_bit = frame_status_bit();
    stream_software_dds_sample();
}
\end{lstlisting}

Interactive UART commands adjust enable, mute, sample width, sample rate, and tone frequency at runtime.
Parameter clamping in software enforces safe DDS limits before writing \texttt{CONTROL}.

During hardware bring-up, three firmware bugs were identified and fixed without RTL changes: (i)~a bit mask that corrupted sample sign, (ii)~incorrect 32-bit sample justification, and (iii)~inaccurate \texttt{usleep()}-based pacing replaced by STATUS polling as shown in \Cref{lst:firmware}.

\section{Experimental Results}
\label{sec:results}

Validation used oscilloscope measurements, listening tests on the PCM5102A, Vivado synthesis reports, and supplementary RTL simulation.
Waveform captures confirm 64 BCLK periods per stereo frame, WS transitions one BCLK before channel MSBs, and a zero delay bit after each WS edge as required by the Philips specification~\cite{philips_i2s_1986}.

\subsection{Arbitrary Sample Rate Accuracy}

Programming \texttt{SAMPLE\_RATE\_HZ} to standard and non-standard values yields the results in \Cref{tab:fs-accuracy}.
Measured WS frequency tracks the programmed target with relative error below 0.001\%, demonstrating the effectiveness of the 32-bit DDS accumulator at 100\,MHz.

\begin{figure}[!t]
  \caption{Measured sample-rate accuracy of DDS-based WS generation.}
  \label{tab:fs-accuracy}
  \centering
  \resizebox{\linewidth}{!}{%
  \begin{tabular}{@{}rrrr@{}}
    \toprule
    \textbf{Target (Hz)} & \textbf{Measured (Hz)} & \textbf{Err. (Hz)} & \textbf{Err. (\%)} \\
    \midrule
    10{,}000  & 9{,}999.98  & $-$0.02 & $-$0.0002 \\
    12{,}345  & 12{,}344.95 & $-$0.05 & $-$0.0004 \\
    44{,}100  & 44{,}099.92 & $-$0.08 & $-$0.0002 \\
    48{,}000  & 47{,}999.95 & $-$0.05 & $-$0.0001 \\
    96{,}000  & 95{,}999.91 & $-$0.09 & $-$0.0001 \\
    \bottomrule
  \end{tabular}%
  }
\end{figure}

Dynamic sweeps from 8\,kHz to 96\,kHz produced clean audible output without perceptible pitch distortion, indicating smooth rate changes through register writes alone.

\subsection{Audio and Control Verification}

Sine tones at 1\,kHz, 2\,kHz, and 5\,kHz matched oscilloscope frequency readouts after the firmware timing fix.
Enable and mute bits in \texttt{CONTROL} correctly gate analog output: \texttt{ENABLE=0} forces lines low, while \texttt{MUTE=1} transmits zero-valued samples with clocks still running.

\begin{figure}[!t]
  \centering
  \IfFileExists{1khz.jpeg}{%
    \includegraphics[width=\linewidth]{1khz.jpeg}%
  }{}
  \caption{Oscilloscope capture of 1\,kHz sine output after firmware correction.}
  \label{fig:1khz}
\end{figure}

\subsection{Resource Utilization}

On the XC7A35T device, the complete SoC including the custom IP uses 2363 LUTs (11.4\%), 2412 flip-flops (5.8\%), and 32 BRAM tiles (64\%).
The modest logic footprint leaves headroom for additional peripherals on the same board.

\section{Conclusion and Future Work}
\label{sec:conclusion}

This paper presented a custom AXI4-Lite I2S transmitter for a MicroBlaze~V SoC on Basys3.
The IP generates Philips-format stereo audio with DDS-based arbitrary sample rates, shadow-register sample integrity, and STATUS-driven firmware handshake.
Hardware measurements confirm correct I2S timing and sub-0.001\% sample-rate accuracy; DAC listening tests verify functional audio output.

Future extensions include TX FIFOs with watermark interrupts, DMA-based streaming for long playback, optional RX or TDM modes, and Linux/ASoC integration.
For applications requiring sub-100\,ppm clock accuracy, an external audio PLL may complement the digital DDS approach.

\section*{Acknowledgment}

The authors thank Ir.\ Agus Bejo, S.T., M.Eng., D.Eng., IPM.\ and Ir.\ Wahyu Dewanto, M.T.\ for supervising the undergraduate thesis from which this paper is derived.


\balance
\bibliographystyle{IEEEtranN}
\bibliography{paper}

\end{document}
